\chapter*{ABSTRAK}
\addcontentsline{toc}{chapter}{ABSTRAK}

\begin{center}

    \textbf{Perancangan Arsitektur Sistem Inspeksi Kargo Digital Terintegrasi untuk Optimalisasi Proses Pemeriksaan Kontainer di Pelabuhan} \\
    \vspace{0.4cm}
    oleh \\
    \vspace{0.4cm}
    Aththariq Lisan Quran Daulah Sentono\\
    18222013\\
\end{center}


\begin{spacing}{1.0}
Proses pemeriksaan kontainer di pelabuhan masih menghadapi berbagai kendala,
di antaranya fragmentasi informasi antarpemangku kepentingan, pencatatan hasil
inspeksi yang belum terintegrasi, keterbatasan ketertelusuran bukti inspeksi,
dan rendahnya visibilitas hasil pemeriksaan secara waktu nyata. Kondisi
tersebut menyebabkan proses verifikasi kontainer, validasi manifes, dan
pengambilan keputusan operasional menjadi kurang efisien. Tugas akhir ini
bertujuan merancang arsitektur sistem inspeksi kargo digital terintegrasi yang
mendukung identifikasi kontainer, pengolahan hasil inspeksi, integrasi data,
dan penyajian informasi secara terpadu.

Perancangan dilakukan dengan pendekatan \textit{Design Science Research} yang
mencakup identifikasi masalah, penetapan tujuan, perancangan artefak,
demonstrasi melalui realisasi sistem, dan evaluasi hasil. Arsitektur sistem
disusun dengan pendekatan berlapis yang mencakup lapisan presentasi, aplikasi,
data, dan \textit{edge}. Pada rancangan ini, komponen \textit{edge} bertugas
melakukan akuisisi aliran video dan inferensi awal, layanan API menangani
orkestrasi serta persistensi data, sedangkan aplikasi web menyediakan antarmuka
pemantauan hasil inspeksi. Pada aspek data, entitas kontainer ditempatkan
sebagai pusat integrasi hasil inspeksi, manifes, dan bukti visual.

Hasil realisasi sistem menunjukkan bahwa rancangan arsitektur dapat diwujudkan
menjadi sistem yang menghubungkan komponen \textit{edge}, layanan API,
penyimpanan data, dan antarmuka web dalam satu alur yang konsisten. Evaluasi
dilakukan melalui pengujian integrasi terhadap alur utama, meliputi pembentukan
entitas kontainer melalui OCR, propagasi nomor kontainer ke alur pemindaian
lain, penyimpanan hasil inspeksi pada model data inti, dan penyajian hasil pada
aplikasi web. Hasil evaluasi menunjukkan bahwa rancangan yang diusulkan mampu
mendukung proses pemeriksaan kontainer yang lebih terstruktur, terdokumentasi,
dan siap dikembangkan lebih lanjut. Dengan demikian, kontribusi utama tugas
akhir ini adalah rancangan arsitektur yang memetakan pembagian tanggung jawab
antarkomponen, menempatkan entitas kontainer sebagai pusat integrasi data
inspeksi, dan memvalidasi kesesuaian rancangan terhadap realisasi sistem.

Kata kunci: arsitektur sistem, inspeksi kargo digital, integrasi data, pelabuhan, ketertelusuran.
\end{spacing}
